
\medskip
\phantomsection\label{ex:sga3-VIA-1.3.2}%
\noindent\textbf{Examples 1.3.2.}
Let \(k\) be a nonperfect field of characteristic \(p>0\), let
\(t\in k-k^p\), let \(\overline k\) be an algebraic closure of \(k\), and
let \(\alpha\in\overline k\) satisfy \(\alpha^p=t\).

\begin{enumerate}
\item[(1)] Consider the additive group
\(\mathbb G_{a,k}=\operatorname{Spec}k[X]\), and let \(G\) be the group
subscheme, finite over \(k\), defined by the additive polynomial
\(X^{p^2}-tX^p\). Then
\[
  G_{\mathrm{red}}
    =\operatorname{Spec}k[X]/X\bigl(X^{p(p-1)}-t\bigr)
\]
is \'etale at the origin. If it were a group scheme over \(k\), it would
be smooth (by \hyperref[prop:sga3-VIA-1.3.1]{1.3.1}); but \(G_{\mathrm{red}}\) is not geometrically
reduced, so it is not a group scheme over \(k\).

\item[(2)] Consider
\[
  \mathbb G^4_{a,k}=\operatorname{Spec}k[X,Y,U,V],
\]
and let \(G\) be the group subscheme defined by the ideal \(I\) generated
by the additive polynomials
\[
  P=X^p-tY^p
  \qquad\text{and}\qquad
  Q=U^p-tV^p.
\]
Then \(G\) is two-dimensional and irreducible, since
\[
  (G_{\overline k})_{\mathrm{red}}
    \simeq\operatorname{Spec}\overline k[Y,V]
\]
is irreducible.

Let \(A=k[X,Y,U,V]\), and let \(\mathfrak m\) be its augmentation ideal.
Denote by \(x,y,u,v\) the images of \(dX,dY,dU,dV\) in
\[
  \Omega^1_{A/k}\otimes_A(A/\mathfrak m),
\]
viewed as linear forms on the tangent space
\[
  k^4=T_0\mathbb G^4_{a,k}.
\]
We shall show that the subspace \(E=T_0G_{\mathrm{red}}\) equals \(k^4\).
Otherwise, there would be a linear form
\[
  f=ax+by+a'u+b'v,
  \qquad a,b,a',b'\in k,
\]
not all of whose coefficients are zero, that vanishes on \(E\). Recall
that the formation of \(\Omega^1_{A/k}\), and hence that of tangent spaces,
commutes with base change (cf.\ EGA~IV\textsubscript{4}, 16.4.5), and
identify \(f\) with its image in \((\overline k^4)^*\). Since
\[
  (G_{\overline k})_{\mathrm{red}}
    \subset (G_{\mathrm{red}})_{\overline k},
\]
the form \(f\) vanishes on the subspace
\(T_0(G_{\overline k})_{\mathrm{red}}\) of \(\overline k^4\), which is
defined by the equations
\[
  g_1=x-\alpha y
  \qquad\text{and}\qquad
  g_2=u-\alpha v.
\]
Thus
\[
  f=\lambda g_1+\mu g_2,
  \qquad \lambda,\mu\in\overline k.
\]
But \(\lambda g_1+\mu g_2\) belongs to \(k^4\) only if
\(\lambda=\mu=0\). This contradiction shows that \(E=k^4\), and hence
\[
  T_0(G_{\mathrm{red}})_{\overline k}=\overline k^4.
\]

On the other hand, \(R=XV-YU\) belongs to \(\sqrt I\), since
\[
  R^p=(X^p-tY^p)V^p-Y^p(U^p-tV^p).
\]
Consequently, the tangent space \(F\) at the point
\((\alpha,1,\alpha,1)\) of \((G_{\mathrm{red}})_{\overline k}\) is
contained in the hyperplane \(H\) of \(\overline k^4\) with equation
\[
  \alpha\,dV+dX-dU-\alpha\,dY=0,
\]
and therefore has dimension at most \(3\).%
\footnote{\label{note:sga3-VIA-16}Editor's note: in fact, it is easy to
see that \(\sqrt I\) is generated by \(P,Q,R\) at every point other than
\(0\), so \(F=H\).}
It follows from \hyperref[concl:sga3-VIA-1.3.1-1]{part~(1)} of \SGARefLink{prop:sga3-VIA-1.3.1}{1.3.1} that \(G_{\mathrm{red}}\) is not a group
scheme over \(k\).
\end{enumerate}

\section*{2. Connected components of an \(A\)-group locally of finite type}
\phantomsection\label{sec:sga3-VIA-2}
\addcontentsline{toc}{section}{2. Connected components of an
  \texorpdfstring{\(A\)}{A}-group locally of finite type}

\medskip
\phantomsection\label{par:sga3-VIA-2.1}%
\noindent\textbf{2.1.}
Let us first consider an arbitrary \(A\)-group \(G\), and let \(G'\) be
the connected component of the origin \(e\) in \(G\). This connected
component is plainly closed, so we may identify it with the reduced closed
subscheme of \(G\) having \(G'\) as its underlying space.%
\footnote{\label{note:sga3-VIA-17}Editor's note: we shall see below (\SGARefLink{par:sga3-VIA-2.2}{2.2})
that \(G'\) is a group in the category
\((\mathbf{Sch}_{/k})_{\mathrm{red}}\).}

\medskip
\phantomsection\label{prop:sga3-VIA-2.1.1}%
\noindent\textbf{Proposition 2.1.1.}
\textit{For every extension \(K\) of the residue field \(k\) of \(A\),
the underlying space of \(G'\otimes_AK\) is the connected component of the
origin in the \(K\)-group \(G\otimes_AK\); in other words, \(G'\) is
geometrically connected.}

\medskip
Indeed, let \((G\otimes_AK)'\) be the connected component of the origin in
\(G\otimes_AK\). Since the image of \((G\otimes_AK)'\) in \(G\) is
connected and contains the unit element of \(G\), that image is contained
in \(G'\). Thus \((G\otimes_AK)'\) is contained in the inverse image
\(G'\otimes_AK\) of \(G'\) in \(G\otimes_AK\). The proposition therefore
follows from the connectedness of \(G'\otimes_AK\), which is proved in the
following lemma.

\medskip
\phantomsection\label{lem:sga3-VIA-2.1.2}%
\noindent\textbf{Lemma 2.1.2.}
\textit{Let \(X\) and \(Y\) be two connected schemes over a field \(k\).
If \(X\) contains a rational point, then \(X\times_kY\) is connected.}

\medskip
We give a direct proof below of this result from EGA~IV\textsubscript{2}
(4.5.8 and 4.5.14).

Suppose first that \(Y\) is nonempty, connected, and affine, with algebra
\(B\). Then \(X\times_kY\) is the spectrum of the quasi-coherent
\(\mathcal O_X\)-algebra
\[
  \mathcal B=\mathcal O_X\otimes_kB.
\]
We wish to show that every nonempty subset \(U\) of \(X\times_kY\) that is
both open and closed equals \(X\times_kY\). Now \(U\) is affine over \(X\),
and its affine \(\mathcal O_X\)-algebra is a direct factor of
\(\mathcal B\). It follows from \hyperref[lem:sga3-VIA-2.1.3]{Lemma~2.1.3} below that the image of \(U\)
in \(X\) is open and closed, and hence is all of \(X\). In particular,
this image contains a rational point \(x\) of \(X\), so \(U\) meets the
inverse image of \(x\) in \(X\times_kY\). Since this inverse image is
isomorphic to \(Y\), and hence connected, \(U\) contains the whole inverse
image. The same conclusion would hold for the complement of \(U\) in
\(X\times_kY\) if \(U\) were different from \(X\times_kY\), which is
absurd.

If \(Y\) is now an arbitrary \(k\)-scheme, the preceding argument shows
that the fibers of the canonical projection
\[
  X\times_kY\longrightarrow Y
\]
are connected. If \(x\) is a rational point of \(X\), all these fibers
meet the subscheme \(\{x\}\times_kY\), which is itself connected. This
proves the assertion.

\medskip
\phantomsection\label{lem:sga3-VIA-2.1.3}%
\noindent\textbf{Lemma 2.1.3.}
\textit{Let \(X\) be a scheme, and let \(\mathcal A\) be a
quasi-coherent \(\mathcal O_X\)-algebra that is locally%
\footnote{\label{note:sga3-VIA-18}Editor's note: the word ``locally'' has
been added. The reference to Kaplansky is \emph{Projective modules},
Ann.\ of Math.\ \textbf{68} (1958), 372--377; see also
\textup{[BAC]}, \S~II.3, Ex.\ 3.}
a direct factor of a free \(\mathcal O_X\)-module. Then the image of
\(\operatorname{Spec}\mathcal A\) in \(X\) is open and closed.}

\medskip
Let \(V\) be this image. It is clear that \(V\) is contained in the
support of \(\mathcal A\). Conversely, if \(x\) belongs to the support of
\(\mathcal A\), then \(\mathcal A_x\) is nonzero and is a free
\(\mathcal O_{X,x}\)-module, since, by Kaplansky, every projective module
over a local ring is free. Consequently, the fiber of
\(\operatorname{Spec}\mathcal A\) at \(x\), which is affine with algebra
\[
  \mathcal A_x\otimes_{\mathcal O_{X,x}}\kappa(x),
\]
is nonempty. Thus the image of \(\operatorname{Spec}\mathcal A\)
coincides with the support of \(\mathcal A\).

If \(s\) is the unit section of \(\mathcal A\), the equality \(s_x=0\)
implies that \(s\), and hence \(\mathcal A\), vanish in a neighborhood of
\(x\). Therefore the support of \(\mathcal A\) is closed. On the other
hand, we may assume that \(\mathcal A\) is a direct factor of the direct
sum \(\mathcal L\) of a family
\((\mathcal O_X^\alpha)\) of copies of \(\mathcal O_X\).%
\footnote{\label{note:sga3-VIA-19}Editor's note: the original has been
expanded in what follows.}
For each \(\beta\), let \(\varphi^\beta\) denote the restriction to
\(\mathcal A\) of the canonical projection
\[
  \mathcal L=\bigoplus_\alpha\mathcal O_X^\alpha
    \longrightarrow\mathcal O_X^\beta.
\]
For \(x\in X\), let \(\mathcal A_x^*\) denote the dual module
\[
  \operatorname{Hom}_{\mathcal O_{X,x}}
    (\mathcal A_x,\mathcal O_{X,x}).
\]
Since \(\mathcal A_x\) is a direct factor of \(\mathcal L_x\), the
restriction map
\[
  \operatorname{Hom}_{\mathcal O_{X,x}}
    (\mathcal L_x,\mathcal O_{X,x})
  \longrightarrow
  \operatorname{Hom}_{\mathcal O_{X,x}}
    (\mathcal A_x,\mathcal O_{X,x})
  =\mathcal A_x^*
\]
is surjective. If \(\mathcal A_x\ne0\), then, since \(\mathcal A_x\) is
free and nonzero, there exist \(a\in\mathcal A_x\) and a linear form
\(\varphi\in\mathcal A_x^*\) such that \(\varphi(a)=1\). There is
therefore a family \((\xi^\alpha)\) of elements of
\(\mathcal O_{X,x}\) such that, for every
\(u\in\mathcal A_{X,x}\),
\[
  \varphi(u)=\sum_\alpha\xi^\alpha\varphi^\alpha(u);
\]
this sum is finite because \(\varphi^\alpha(u)=0\) except for finitely many
\(\alpha\). Applying this to \(u=a\), we find
\(\alpha_1,\ldots,\alpha_n\) such that
\[
  1=\varphi(a)
   =\xi^{\alpha_1}\varphi^{\alpha_1}(a)
    +\cdots+
    \xi^{\alpha_n}\varphi^{\alpha_n}(a).
\]
There is then an open neighborhood \(U\) of \(x\) on which \(a\) and the
\(\xi^{\alpha_i}\) come from sections
\[
  \widetilde a\in\mathcal A(U)
  \qquad\text{and}\qquad
  \widetilde\xi^{\alpha_i}\in\mathcal O_X(U).
\]
The equality
\[
  \sum_i
    \widetilde\xi^{\alpha_i}
    \varphi^{\alpha_i}(\widetilde a)=1
  \quad\text{on }U
\]
shows that \(\widetilde a_y\ne0\) for every \(y\in U\).
(``The support of a projective module is open.'')

\medskip
\phantomsection\label{par:sga3-VIA-2.2}%
\noindent\textbf{2.2.}
The notation remains that of \hyperref[par:sga3-VIA-2.1]{2.1}. It is clear that \(G'\) is a reduced
\(k\)-scheme. \hyperref[lem:sga3-VIA-2.1.2]{Lemma~2.1.2} shows that \(G'\times_kG'\) is connected, so
\((G'\times_kG')_{\mathrm{red}}\) is the reduced subscheme of
\(G\times_AG\) whose underlying space is the connected component of the
origin. In particular, the multiplication morphism
\[
  \mu:G\times_AG\longrightarrow G
\]
induces a morphism
\[
  \mu':(G'\times_kG')_{\mathrm{red}}\longrightarrow G'
\]
that makes \(G'\) a group in
\((\mathbf{Sch}_{/k})_{\mathrm{red}}\).

\medskip
\phantomsection\label{par:sga3-VIA-2.2bis}%
\noindent\textbf{2.2.bis.}%
\footnote{\label{note:sga3-VIA-20}Editor's note: this paragraph has been
added in order to make the connection with
Expos\'e~VI\textsubscript{B}, \S~3.}
Recall (cf.\ Expos\'e~V) that, if \(P\) is a scheme, \(\underline P\)
denotes the underlying topological space of \(P\). We define an
\(A\)-subfunctor \(G^0\) of \(G\) by setting, for every \(A\)-scheme
\(S\),
\[
  G^0(S)
    =\{u\in G(S)\mid u(\underline S)\subset\underline{G'}\}.
\]
Let \(c:G\to G\) be the inversion morphism. Since
\(c(\underline{G'})=\underline{G'}\), we have
\(c\circ u\in G^0(S)\) for every \(u\in G^0(S)\). On the other hand, if
\(u,v\in G^0(S)\), then \(u\boxtimes v\) sends \(\underline S\) into the
subspace of \(\underline{G\times_AG}\) consisting of the points whose two
projections belong to \(\underline{G'}\). This subspace is identified with
the underlying space of \(G'\times_AG'\), which is connected by
\hyperref[lem:sga3-VIA-2.1.2]{Lemma~2.1.2}. Consequently, \(\mu\circ(u\boxtimes v)\) sends
\(\underline S\) into \(\underline{G'}\). This proves that \(G^0\) is a
group subfunctor of \(G\) over \(A\).

If the connected component of \(e\) is open in \(G\), then the subfunctor
\(G^0\) is represented by the subscheme induced by \(G\) on this open
subset, which is therefore a group subscheme of \(G\); this subscheme will
also be denoted by \(G^0\). In this case, with the notation of \SGARefLink{prop:sga3-VIB-2.1}{2.1},
\[
  G'=(G^0)_{\mathrm{red}},
\]
and the topological spaces \(\underline{G'}\) and \(\underline{G^0}\)
coincide.%
\footnote{\label{note:sga3-VIA-21}Editor's note: in this expos\'e, the
notation \(G^0\) is reserved for the case in which the connected component
of \(e\) is open; in Expos\'e~VI\textsubscript{B}, \S~3, this connected
component will be denoted by \(\underline{G^0}\) in every case. This is a
slight abuse of notation, but it is compatible with the preceding
discussion when the connected component is the underlying topological
space of an open group subscheme \(G^0\) of \(G\).}

\medskip
\phantomsection\label{par:sga3-VIA-2.3}%
\noindent\textbf{2.3.}
In accordance with our convention in \SGARefLink{par:sga3-VIA-1.1}{1.1}, we again assume from now on that
\(G\) is locally of finite type over \(A\). Then \(G\) is locally
Noetherian, hence locally connected;%
\footnote{\label{note:sga3-VIA-22}Editor's note: indeed, in a Noetherian
space there are only finitely many connected components, and each is
therefore open; see also EGA~I, 6.1.9.}
thus every connected component of \(G\) is open.

We shall write \(G^0\) for the subscheme induced by \(G\) on the connected
component \(G'\) of the unit element. By \hyperref[par:sga3-VIA-2.2bis]{2.2.bis}, \(G^0\) is a group
subscheme of \(G\), which we shall call the \emph{neutral component} of
\(G\). For every \(A\)-scheme \(S\), therefore,
\[
  G^0(S)
    =\{u\in G(S)\mid
      u(\underline S)\subset\underline{G^0}=\underline{G'}\}.
\]

Let \(G^\alpha\) be an arbitrary connected component of \(G\), and let
\[
  \nu^\alpha:G^\alpha\times_AG^0\longrightarrow G
\]
be the morphism defined by
\[
  \nu^\alpha(S)(g,\gamma)=g\gamma g^{-1}
\]
for every \(S\in(\mathbf{Sch}_{/A})\), \(g\in G^\alpha(S)\), and
\(\gamma\in G^0(S)\). If \(e\) is the origin of \(G\), the restriction of
\(\nu^\alpha\) to \(G^\alpha\times_A\{e\}\) is the trivial morphism.
Since \(G^\alpha\times_AG^0\) is connected by \hyperref[lem:sga3-VIA-2.1.2]{2.1.2}, it follows that
\(\nu^\alpha\) factors through \(G^0\). Thus, for every \(A\)-scheme
\(S\), \(G^0(S)\) is a normal subgroup of \(G(S)\). We have proved the
following proposition.%
\footnote{\label{note:sga3-VIA-23}Editor's note: the number \SGARefLink{prop:sga3-VIA-2.3.1}{2.3.1} has
been added in order to emphasize this statement. Moreover, \(G^0\) is
even a characteristic subgroup of \(G\); see \SGARefLink{thm:sga3-VIA-2.6.5:item:ii}{2.6.5(ii)}.}

\medskip
\phantomsection\label{prop:sga3-VIA-2.3.1}%
\noindent\textbf{Proposition 2.3.1.}
\textit{Let \(G\) be an \(A\)-group locally of finite type. Then the
neutral component \(G^0\) is an open normal group subscheme of \(G\).}

\medskip
\phantomsection\label{prop:sga3-VIA-2.4}%
\noindent\textbf{Proposition 2.4.}
\textit{Let \(G\) be an \(A\)-group locally of finite type.}
\begin{enumerate}
\item[\SGAInlineRefTarget{prop:sga3-VIA-2.4:item:i}(i)] \textit{\(G^0\) is irreducible, and \(G^0\otimes_Ak\) is
geometrically irreducible over \(k\).}
\item[\SGAInlineRefTarget{prop:sga3-VIA-2.4:item:ii}(ii)] \textit{\(G^0\) is quasi-compact, and hence is of finite type
over \(A\).}
\end{enumerate}

\noindent\footnote{\label{note:sga3-VIA-24}Editor's note: in the
statement, ``\(G^0\) is geometrically irreducible'' has been replaced by
``\(G^0\otimes_Ak\) is geometrically irreducible over \(k\),'' and the
proof has been expanded.}
\textit{Proof.}
\SGARefLink{prop:sga3-VIA-2.4:item:i}{(i)} Since \(G^0\) and \(G^0\otimes_Ak\) have the same underlying
topological space, it suffices to prove the second assertion. Let
\(\overline k\) be an algebraic closure of \(k\). By \SGARefLink{prop:sga3-VIB-2.2}{2.2},
\((G\otimes_A\overline k)_{\mathrm{red}}\) is a reduced
\(\overline k\)-group locally of finite type, and hence is smooth over
\(\overline k\) by \SGARefLink{prop:sga3-VIA-1.3.1}{1.3.1}. A fortiori, the local rings of
\((G\otimes_A\overline k)_{\mathrm{red}}\) are integral domains.
Therefore,%
\footnote{\label{note:sga3-VIA-25}Editor's note: the original has been
expanded here, and the reference to EGA~I, 6.1.10, has been added.}
since \(G\otimes_A\overline k\) is locally Noetherian, the connected
components of \(G\otimes_A\overline k\) are irreducible
(cf.\ EGA~I, 6.1.10). In particular, the connected component
\(G^0\otimes_A\overline k\) (cf.\ 2.1.1) is irreducible.

\SGARefLink{prop:sga3-VIA-2.4:item:ii}{(ii)} We now show that \(G^0\) is of finite type over \(A\). Since \(G^0\)
is locally of finite type over \(A\), it suffices to prove that \(G^0\) is
quasi-compact. Since \(G^0\) is irreducible, this follows from \SGARefLink{sga3:VIA:0.5.1}{0.5.1}.

\medskip
\phantomsection\label{cor:sga3-VIA-2.4.1}%
\noindent\textbf{Corollary 2.4.1.}
\textit{Every connected component of \(G\) is irreducible,%
\footnote{\label{note:sga3-VIA-26}Editor's note: one should bear in mind
that a nonneutral connected component is not geometrically connected in
general. For example, if \(k=\mathbb R\), the group
\(\mu_{3,\mathbb R}\), represented by
\(\mathbb R[X]/(X^3-1)\), has two connected components:
\(\{e\}=\operatorname{Spec}\mathbb R\) and
\(C=\operatorname{Spec}\mathbb R[X]/(X^2+X+1)\); and
\(C\otimes_{\mathbb R}\mathbb C\) has two components.}
of finite type over \(A\),%
\footnote{\label{note:sga3-VIA-27}Editor's note: the assertion that
follows has been added; cf.\ Expos\'e~VI\textsubscript{B}, \SGARefLink{lem:sga3-VIB-1.5}{1.5}.}
and of the same dimension as \(G^0\).}

\medskip
Indeed, we may assume that \(A\) equals its residue field \(k\). Let \(C\)
be a connected component of \(G\), let \(x\) be a closed point of \(C\),
let \(\kappa(x)\) be the residue field of \(x\), and let \(k'\) be a
finite normal extension of \(k\) containing \(\kappa(x)\). The canonical
projection
\[
  \pi:C\otimes_kk'\longrightarrow C
\]
is open and closed.%
\footnote{\label{note:sga3-VIA-28}Editor's note: the original has been
expanded in what follows.}
Consequently, if \(C'\) is a connected component of \(C\otimes_kk'\), the
projection \(C'\to C\) is surjective. Thus \(C'\) contains a point
\(y\in\pi^{-1}(x)\), and such a point is rational over \(k'\)
(cf.\ the proof of \hyperref[prop:sga3-VIA-1.2]{1.2}); it follows that \(C'\) is the image of
\(G^0\otimes_kk'\) under right translation \(r_y\). Now
\(G^0\otimes_kk'\) is of finite type over \(k'\), by \hyperref[prop:sga3-VIA-2.4]{2.4}, and
\(\pi^{-1}(x)\) is finite, of cardinality at most \([k':k]\).
Consequently, \(C\otimes_kk'\) is of finite type over \(k'\), and hence
\(C\) is of finite type over \(k\).

Moreover, since \(G^0\otimes_kk'\) is irreducible by \SGARefLink{prop:sga3-VIA-2.4}{2.4}, the same is true
of \(C'\),%
\footnote{\label{note:sga3-VIA-29}Editor's note: the original has been
simplified here.}
and hence also of \(C\), because the projection \(C'\to C\) is
surjective.

Finally, we saw above that \(C\otimes_kk'\) is the disjoint union of
finitely many translates of \(G^0\otimes_kk'\). Since dimension is
invariant under extension of the base field (cf.\
EGA~IV\textsubscript{2}, 4.1.4), it follows that \(C\) has the same
dimension as \(G^0\). (Moreover, by EGA~IV\textsubscript{2}, 5.2.1, one
has \(\dim_gG=\dim G^0\) for every point \(g\in G\).)
